% Methodology additions -- drop-in fragment for the paper.
% Requires: \usepackage{booktabs}. Slots after the existing Section 2.
% These subsections extend the draft's Methodology to cover the experiments
% that were actually run (dial sweep, native cross-language collection, and the
% garbage-collector ablation) and record the dataset-governance rules.

\subsection{Dial Sweep at a Fixed Trace Count}
\label{sec:method-dial}
The threshold experiment varies trace count while holding the two collection
``dials'' fixed. A complementary experiment holds trace count fixed and varies
the dials, to separate the contribution of \emph{averaging} from that of
\emph{signal strength}:
\begin{itemize}
  \item \texttt{-repeat} $r$: the number of vulnerable encryptions accumulated
        into one trace. Larger $r$ raises the target operation's contribution
        relative to fixed timer overhead (more averaging per trace).
  \item \texttt{-evict-kb} $e$: the size of the memory region read before each
        timed operation. Larger $e$ disturbs more of the cache and is expected
        to strengthen the collision signal.
\end{itemize}
The count is fixed at \textbf{100{,}000 traces}. This value is chosen
deliberately: it sits in the C attack's \emph{marginal} region, just above its
50\% threshold (Section~\ref{sec:results-threshold}), where recovery is most
sensitive to changes in signal quality. A count deep in the reliable regime
(where recovery is already~100\%) or deep in the failure regime (0\%) would
mask any dial effect; the knee maximises the visible dynamic range. We sweep
$r \in \{10,20,30,50,75,100\}$ and $e \in \{256,512,1024,2048,4096\}$~KiB, each
about the $r=50$, $e=2048$ baseline, with ten interleaved trials per point.

\subsection{Native Cross-Language Recovery}
\label{sec:method-crosslang}
The draft's cross-language plan (Section~2.11) supplies \emph{one shared trace
file} to C, Go, and Python to test algorithmic consistency and processing time.
The experiments reported here answer a different question. Each implementation
performs its \emph{own} measured collection with its native timer and then runs
its own recovery, and we compare the resulting \emph{recovery reliability} as a
function of trace count. This measures how each language's runtime affects the
quality of the timing channel it can capture---the ``separate runtime-dependent
experiment'' the draft anticipates---rather than the correctness of a shared
computation.

Because each language collects independently, the trace data (and its random
seed) differ across implementations. Trace count, the fixed AES key, the dial
settings, the outlier rule, the recovery algorithm, and the verification
criterion are held identical; only the implementation (and hence its collection
runtime) changes. Differences in recovery are therefore attributable to the
language's runtime behaviour during collection, not to the attack logic, which
is common. We treat this as a \emph{native-collection} comparison and keep it
distinct from the shared-dataset consistency test, which remains future work.

\subsection{Garbage-Collector Ablation (Go)}
\label{sec:method-gc}
Go executes within a managed runtime whose garbage collector (GC) and goroutine
scheduler can preempt execution at unpredictable moments, adding jitter to the
per-trace timings the attack depends on. To test whether this runtime noise is
the cause of Go's higher threshold, we repeat the Go collection with two
environment settings that quiet the runtime without altering the attack:
\texttt{GOGC=off} (disables the collector, so no GC pause lands inside a timed
region) and \texttt{GOMAXPROCS=1} (pins execution to one core, suppressing
cross-core migration). Both are read by the Go runtime from the environment; no
source change is made. For cost reasons the ablation uses a reduced grid (counts
$\{500{,}000, 250{,}000, 100{,}000, 50{,}000\}$ and a $3\times3$ dial grid), so
it is compared to the garbage-collected run only at matched points and by the
fitted 50\% threshold, never by a grid-dependent high-reliability count.

\subsection{Dataset Governance and Exclusion Criteria}
\label{sec:method-governance}
Several exploratory runs preceded the retained runs. To avoid selective
reporting, inclusion follows fixed rules, and every excluded run's metadata is
preserved (see the provenance table).
\begin{itemize}
  \item \textbf{Environment gate.} A run is excluded if its mean collection cost
        exceeds roughly $1.2$~ms/sample (baseline $\approx 0.94$), a proxy for
        background load. One 10-trial C sweep is excluded on this basis
        ($\approx 3.0$~ms/sample, threefold inflated).
  \item \textbf{Design gate.} Pilot runs with fewer trials, missing
        reproducibility sidecars, or single-point grids are excluded from the
        primary analysis and used only for design context.
  \item \textbf{Implementation identity.} One retained C run was launched with an
        \texttt{IMPL=go} tag before the harness supported language selection;
        the C binary in fact executed. It is verified as C (no \texttt{impl}
        field in its metadata; \texttt{min\_time} and ms/sample matching C) and
        relabelled accordingly. This mislabelling is documented because an
        earlier positional (rather than by-name) parse of the result column
        produced incorrect tallies; all reported counts are recomputed by column
        name.
\end{itemize}
